% s-020-p020-m002; original PDF 20, 21
\Aphorism{s-020-p020-m002}{10}
Kant’s Denken ist so wurzelhaft deterministisch, dass ihm die tiefste sittliche Ekstase in einem „Du sollst“, die klarste Erkenntniss der Erfahrungswelt im starren Dodekalog einer Kategorientafel aufging; wo er das Letzte seiner Seele blosslegt, ist immer Pflicht, Zwang, Formalismus in der Nähe — er versteht sich und die Dinge erst, wenn er ein Müssen und Gehorchen darin aufgespürt hat. Sollte nicht wirklich, wie man oft halb im Scherz zugestand, der Staat Friedrich’s des Grossen, der preussische Beamtenstaat diesen Denker und seine Idiosynkrasie auf dem Gewissen haben? Sowie der reactionäre Polizeistaat den Denker Hegel, sowie umgekehrt das neue deutsche Reich mit seiner liberalen Lockerung des gesellschaftlichen Lebens einen modernen Denker auf dem Gewissen haben würde, der indeterministisch bis zum Äussersten empfände und dem, entgegengesetzt zu Kant, Causalität erst verständlich wäre, wenn er die dahinterliegende Freiheit aufgedeckt sähe? Sicherlich sind wir, ohne es zu wissen, von politischen und geschichtlichen Zuständen tief beeinflusst, wenn wir bei der philosophischen Deduction unser Axiom, unsere \Foreign{petitio principii}, \Sperr{unser} \Greek{ποῦ} \Greek{στῶ} formuliren; sicherlich sind es nicht vorzugsweise intellectuelle Quellen, aus denen die Grundströmungen und Haupttendenzen unserer intellectuellen Thätigkeit fliessen. Wie sehr wirthschaftliche und politische Gegebenheiten auf den „Zeitgeist“, auf die Gruppenbestrebungen und gemeinsamen Bildungsziele einer Epoche abfärben, darüber wird seit Langem nachgesonnen; wie sehr unterschieden aber auch der einzelne Mensch zu den objectivsten Problemen des Denkens stehen kann, wie präformirt durch Umgebung, Rasse, Klima, Organisation, physiologisches Gemeingefühl, wie wenig selbst in Dingen des „reinen Geistes“ der — reine Geist zu sagen hat, das ist eine moderne Bedenklichkeit und Erwägung, die in der Theorie vom milieu nur einen vorläufigen, groben und doctrinären Ausdruck gefunden hat. So leichten Kaufes wird man das souveräne Ich nicht los \Dots{3}
\par

